\scanpage{37}
\mspage{18}
% archive foliation: 36

$F(\Delta)/\{1,\underline a_\Delta\}$ (cf (36)), qui est donc l'ens des quatre bifaces de $\Delta$. Donc c'est l'octaèdre orienté canoniquement associé à l'octaèdre $\Delta$, savoir
\[
(45)\quad
\left\{\begin{array}{l}
\Delta^+=\deleted{\unsure{\Delta_\Delta\wedge\mathbb{F}_2}}\mathrm{Or}(\Delta)\wedge_{\{1,\underline a_\Delta\}}\Delta\\
\Delta=\mathrm{Or}(\Delta)\wedge_{\{1,\underline a_\Delta\}}\Delta^+
\end{array}\right.
\]
de sorte qu'on a \inserted{bien}, comme il se doit
\[
(\unsure{45|46})\quad
\left\{\begin{array}{l}
\mathrm{Or}(\Delta^+)\simeq\{\pm1\}=1_{\mathbb{F}_2}\\
\underline\omega_{\Delta^+}\simeq\underline\omega_\Delta\simeq\underline\omega_J=\underline\omega
\quad[\text{car }\underline a_\Delta\text{ opère trivialement sur }\underline\omega_\Delta]\\
C_{\Delta^+}\simeq C_\Delta\wedge\underline\mu\simeq\underline\omega\qquad\ldots
\end{array}\right.
\]

$\big[\big[$ Il importe maintenant de bien comprendre la structure d'extension
\[
(47)\qquad
\diagram{\begin{tikzcd}
1\arrow[r]&\underset{\{\pm1\}}{\underset{\simeq}{Z}}\arrow[r]&
GL(M)\arrow[r]&
\underset{\substack{\mathrm{Aut}_{\mathrm{oct}}(\Delta_M)'\\
\simeq\mathrm{Aut}_{\mathrm{oct}}^+(\Delta_M)\times\{1,\underline a_{\Delta_M}\}}}
{\underset{\simeq}{GL(M)'}}\arrow[r,dash]&1
\end{tikzcd}}
\]
et en particulier
\[
(48)\qquad
\diagram{\begin{tikzcd}
1\arrow[r]&\underset{\{\pm1\}}{\underset{=}{Z_0}}\arrow[r]&
GL(2,\mathbb{Z}/4\mathbb{Z})\arrow[r]&
\underset{\substack{\mathrm{Aut}_{\mathrm{oct}}(\Delta_{(\mathbb{Z}/4\mathbb{Z})^2})\\
\text{``octaèdre standart''}}}
{\underset{\simeq}{GL(2,\mathbb{Z}/4\mathbb{Z})'}}\arrow[r,dash]&1
\end{tikzcd}}
\]
puisque c'est par là qu'on comprend quelle est la donnée supplémentaire qu'il faut pour déduire, d'un octaèdre non orienté $\Delta$, le $(\mathbb{Z}/4\mathbb{Z})$-module libre de rang $2$ $M$ dont $\Delta$ proviendrait — question qui revient à celle de ``remonter'' un torseur sous le groupe $GL(2,\mathbb{Z}/4\mathbb{Z})'$ en un torseur sous le groupe $GL(2,\mathbb{Z}/4\mathbb{Z})$ lui-même.
